% !TEX root = ../main.tex
\chapter{積と余積：レコード、ペア、Either}
\label{chap:products-coproducts}

\begin{chapterabstract}
本章では、圏論の基本構成である\textbf{積}と\textbf{余積}を、ソフトウェアのデータ設計として読む。積は「両方を持つ」構造であり、ペア、レコード、複数フィールドを持つDTOとして現れる。余積は「どちらか一方」の構造であり、\texttt{Either}、\texttt{Result}、エラー分岐、選択肢を持つレスポンス型として現れる。

圏論では、積や余積を単なるデータ構造としてではなく、「その構造を作る・使う関数がどのように一意に決まるか」で特徴づける。この見方を\textbf{普遍性}という。本章では、Lean 4 の \texttt{Prod}、構造体、\texttt{Sum}、\texttt{Option} を使い、積と余積の性質を小さく証明する。最後に、APIレスポンスやエラーハンドリングの設計を題材に、積で表すべき情報と余積で表すべき分岐を区別する。
\end{chapterabstract}

\begin{learninggoals}
\item 積を「両方を持つ」構造として説明できる。
\item 余積を「どちらか一方」の構造として説明できる。
\item \texttt{Prod}、構造体、\texttt{Sum}、\texttt{Option} の対応を読める。
\item 普遍性を「作り方・使い方が一意に決まる」という実務的な言葉で説明できる。
\item エラーを値と同時に持つべきか、成功または失敗の分岐として持つべきかを判断できる。
\end{learninggoals}

\section{この章を読む理由}

APIのレスポンス型を設計するとき、次のような迷いが起きる。

\begin{itemize}
\item ユーザーIDと名前をまとめたい。
\item 成功時には値を返し、失敗時にはエラーを返したい。
\item エラー情報と成功値を同じ構造体に入れてよいのか迷う。
\item \texttt{null}、\texttt{Option}、\texttt{Either}、例外のどれを使うべきか判断したい。
\end{itemize}

これらは、すべて「データをどう組み合わせるか」という問題である。ただし、組み合わせ方には二種類ある。ひとつは、複数の情報を同時に持つ組み合わせである。もうひとつは、複数の候補のうち一つだけを選ぶ組み合わせである。

図\ref{fig:ch13-product-coproduct-overview}では、二つの成分を同時に持つ積と、左右いずれかの注入から作られる余積を対比する。ログインIDとトークンが両方必要なら積、成功値またはエラーの一方だけを返すなら余積として読む。

\FigurePlaceholder{fig-ch13-product-coproduct-overview.png}{0.90}{積と余積}{fig:ch13-product-coproduct-overview}

\section{ソフトウェア上の問題：同時に持つのか、分岐するのか}

\subsection{DTOは多くの場合、積である}

たとえば、次のようなユーザー概要を考える。

\begin{listing}[htbp]
\begin{minted}{lean4}
structure UserSummary where
  id   : Nat
  name : String
\end{minted}
\caption{\texttt{UserSummary}}
\label{lst:ch13_record_intuition}
\end{listing}

この値は、\texttt{id} と \texttt{name} の両方を持つ。どちらか片方ではない。したがって、構造としては積に近い。フィールド名が付いているので、実務コードではペアよりも読みやすいが、圏論的には「複数の射影を持つ値」として理解できる。

\subsection{エラーレスポンスは多くの場合、余積である}

一方、APIレスポンスが次のどちらかである場面を考える。

\begin{itemize}
\item 成功したので、本文を返す。
\item 失敗したので、エラーを返す。
\end{itemize}

この場合、成功値とエラー値が同時にあるとは限らない。むしろ「成功または失敗」のどちらか一方である。ここで使うべき形は、積ではなく余積である。

\texttt{(Error, Value)} は両方を持つ積だが、\texttt{Either Error Value} や \texttt{Result Value Error} はどちらか一方を持つ余積である。この違いを型に反映すると、成功と失敗を同時に有効とみなす状態を排除できる。

\section{積：両方持つ構造}
\label{sec:ch13-product}

積は、二つの型 \(A\) と \(B\) から、両方を同時に持つ型 \(A \times B\) を作る構成である。Leanでは \texttt{Prod A B}、または記法として \texttt{A × B} を使う。

\begin{definitionbox}
型 \(A\) と型 \(B\) の積は、\(A\) の値と \(B\) の値を一つずつ持つ型である。積からは、第一成分を取り出す射影 \(\pi_1 : A \times B \to A\) と、第二成分を取り出す射影 \(\pi_2 : A \times B \to B\) がある。
\end{definitionbox}

Leanでは、ペア \texttt{(a, b)} を作り、\texttt{p.1} と \texttt{p.2} で取り出す。

\begin{listing}[htbp]
\begin{minted}{lean4}
namespace Chapter13

-- Lean の A × B は Prod A B の記法である。
def loginPair : Nat × String := (42, "token")

example : loginPair.1 = 42 := rfl
example : loginPair.2 = "token" := rfl

-- ペアは、取り出した二つの成分から元に戻せる。
theorem prod_eta {A B : Type} (p : A × B) : (p.1, p.2) = p := by
  cases p
  rfl
\end{minted}
\caption{\texttt{Prod} と \texttt{prod\_eta}}
\label{lst:ch13_prod_basic}
\end{listing}

\texttt{(a, b)} はペアを作り、\texttt{p.1} と \texttt{p.2} は二つの射影である。定理 \texttt{prod\_eta} は、ペアから両成分を取り出して作り直すと元のペアに戻ることを示す。

\subsection{レコードは名前付きの積}

実務では、\texttt{Nat × String} のような位置だけで意味を決めるペアよりも、名前付きフィールドを持つ構造体のほうが安全である。

\begin{listing}[htbp]
\begin{minted}{lean4}
structure LoginInput where
  userId : Nat
  token  : String

-- レコードをペアへ変換する。
def loginAsPair (x : LoginInput) : Nat × String :=
  (x.userId, x.token)

-- ペアからレコードを作る。
def pairAsLogin (p : Nat × String) : LoginInput :=
  { userId := p.1, token := p.2 }

theorem pair_login_roundtrip (p : Nat × String) :
    loginAsPair (pairAsLogin p) = p := by
  cases p
  rfl
\end{minted}
\caption{\texttt{LoginInput} と \texttt{Nat × String}}
\label{lst:ch13_record_as_product}
\end{listing}

ペアは軽量だが、フィールドの意味を位置に依存させる。レコードは少し冗長だが、フィールド名によって意図が明確になる。圏論の積としては同じ方向の構造を持つが、実務上の可読性や保守性は異なる。

\begin{warningbox}
積として読めるからといって、常にペアを使えばよいわけではない。業務上の意味がある値には、名前付きフィールドを持つ構造体を使うほうが安全である。圏論は「構造の形」を教えてくれるが、命名やドメイン表現の責任まで自動化してくれるわけではない。
\end{warningbox}

\section{余積：どちらか一方の構造}
\label{sec:ch13-coproduct}

余積は、二つの型 \(A\) と \(B\) から、\(A\) の値か \(B\) の値のどちらか一方を持つ型を作る構成である。Leanでは \texttt{Sum A B} を使う。多くの言語で \texttt{Either A B} と呼ばれる型に対応する。

\begin{definitionbox}
型 \(A\) と型 \(B\) の余積は、\(A\) 側の値または \(B\) 側の値を持つ型である。\(A\) の値を余積へ入れる注入 \(\iota_1 : A \to A + B\) と、\(B\) の値を余積へ入れる注入 \(\iota_2 : B \to A + B\) がある。
\end{definitionbox}

Leanの \texttt{Sum A B} では、左側を \texttt{Sum.inl a}、右側を \texttt{Sum.inr b} として作る。

\begin{listing}[htbp]
\begin{minted}{lean4}
inductive ParseError where
  | empty
  | badFormat

abbrev ParseResult := Sum ParseError Nat

-- 左側はエラー、右側は成功値として読むことにする。
def parseEmpty : ParseResult := Sum.inl ParseError.empty
def parseOk    : ParseResult := Sum.inr 200

-- Sum を使う側は、左右の分岐を両方扱う。
def showParseResult (r : ParseResult) : String :=
  match r with
  | Sum.inl _ => "parse failed"
  | Sum.inr _n => "parse ok"
\end{minted}
\caption{\texttt{ParseResult}}
\label{lst:ch13_sum_basic}
\end{listing}

\texttt{inductive} の各 \texttt{| name ...} は値を作る形を列挙し、\texttt{abbrev} は既存の型に別名を付ける。\texttt{Sum} を受け取る関数は左右両方を処理する必要があるため、成功ケースだけを見てエラーケースを忘れる設計ミスを減らせる。

\section{\texttt{Option} は小さな余積として読める}

\texttt{Option A} は、値がないか、\texttt{A} の値があるかを表す。これは「値なし」という一つだけのケースと、「値あり」というケースの余積として読める。

Leanでは、\texttt{Unit} は値を一つだけ持つ型である。したがって、\texttt{Sum Unit A} は「左なら値なし、右なら値あり」と読め、\texttt{Option A} と同型になる。

\begin{listing}[htbp]
\begin{minted}{lean4}
def optionToSum {A : Type} : Option A → Sum Unit A
  | none   => Sum.inl ()
  | some a => Sum.inr a

def sumToOption {A : Type} : Sum Unit A → Option A
  | Sum.inl _ => none
  | Sum.inr a => some a

theorem option_roundtrip {A : Type} (x : Option A) :
    sumToOption (optionToSum x) = x := by
  cases x with
  | none => rfl
  | some a => rfl

theorem sum_roundtrip {A : Type} (x : Sum Unit A) :
    optionToSum (sumToOption x) = x := by
  cases x with
  | inl u =>
      cases u
      rfl
  | inr a => rfl
\end{minted}
\caption{\texttt{Option A} と \texttt{Sum Unit A}}
\label{lst:ch13_option_sum}
\end{listing}

二つの定理が両方向の round-trip を与える。この見方により、\texttt{Option}、\texttt{Either}、\texttt{Result}、\texttt{Except} を、同じ「分岐を型で表す」設計パターンとして理解できる。

\section{普遍性をソフトウェア的に読む}

積・余積の標準定義は、対象の内部表現ではなく、射影・挿入と媒介射の存在および一意性によって与えられる~\cite{leinster-basic-category-theory,riehl-category-theory-in-context}。
\label{sec:ch13-universal-property}

積と余積は、単に「データの形」として説明できる。しかし圏論では、より重要な説明がある。それが普遍性である。

積の普遍性は、射影 \(\pi_1,\pi_2\) を固定したとき、任意の \(f:C\to A\) と \(g:C\to B\) に対し、\(\pi_1\circ h=f\)、\(\pi_2\circ h=g\) を満たす \(h:C\to A\times B\) がただ一つ存在するという条件である。余積では向きが反転し、二つの注入から出る射を指定すると、余積全体から出る射がただ一つ決まる。ここで「一意」とは、同じ可換条件を満たす候補の間での一意性である。

\subsection{ペアを作る唯一の方法}

同じ入力 \(C\) から \(A\) を作る関数 \(f : C \to A\) と、\(B\) を作る関数 \(g : C \to B\) があるとする。このとき、\(C\) から \(A \times B\) へ行く自然な関数は、\(c\) を \((f(c), g(c))\) に送る関数である。

図\ref{fig:ch13-product-universal}では、\(f\) と \(g\) を二つの射影として再現する媒介射 \(h\) の位置を示す。

\FigurePlaceholder{fig-ch13-product-universal.png}{0.90}{積の普遍性}{fig:ch13-product-universal}

\begin{listing}[htbp]
\begin{minted}{lean4}
def makePair {A B C : Type} (f : C → A) (g : C → B) : C → A × B :=
  fun c => (f c, g c)

theorem makePair_fst {A B C : Type}
    (f : C → A) (g : C → B) (c : C) :
    (makePair f g c).1 = f c := rfl

theorem makePair_snd {A B C : Type}
    (f : C → A) (g : C → B) (c : C) :
    (makePair f g c).2 = g c := rfl
\end{minted}
\caption{\texttt{makePair}}
\label{lst:ch13_product_universal}
\end{listing}

このコードは、積の普遍性のうち「二つの関数からペアを作る」部分を表している。さらに、第一成分と第二成分がそれぞれ \(f\) と \(g\) に一致するなら、そのペア関数は点ごとに一意である。

\begin{listing}[htbp]
\begin{minted}{lean4}
theorem makePair_unique_pointwise {A B C : Type}
    (f : C → A) (g : C → B) (h : C → A × B)
    (hfst : ∀ c, (h c).1 = f c)
    (hsnd : ∀ c, (h c).2 = g c)
    (c : C) : h c = makePair f g c := by
  calc
    h c = ((h c).1, (h c).2) := by
      cases h c
      rfl
    _ = (f c, g c) := by
      rw [hfst c, hsnd c]
    _ = makePair f g c := rfl
\end{minted}
\caption{\texttt{makePair\_unique\_pointwise}}
\label{lst:ch13_product_unique}
\end{listing}

ここでは関数そのものの等しさではなく、任意の入力 \texttt{c} に対する点ごとの等しさとして書いている。関数外延性は第21章であらためて扱うため、本章ではこの形で十分である。

\subsection{分岐を処理する唯一の方法}

余積では向きが逆になる。\(A\) を \(C\) へ処理する関数 \(f : A \to C\) と、\(B\) を \(C\) へ処理する関数 \(g : B \to C\) があるとする。この二つを与えると、\(A+B\) 全体を \(C\) へ処理する関数が一意に決まる。

図\ref{fig:ch13-coproduct-universal}では、左右の注入に続けたとき \(f\) と \(g\) を再現する媒介射を示す。

\FigurePlaceholder{fig-ch13-coproduct-universal.png}{0.90}{余積の普遍性}{fig:ch13-coproduct-universal}

\begin{listing}[htbp]
\begin{minted}{lean4}
def eitherElim {A B C : Type} (f : A → C) (g : B → C)
    (s : Sum A B) : C :=
  match s with
  | Sum.inl a => f a
  | Sum.inr b => g b

theorem eitherElim_inl {A B C : Type}
    (f : A → C) (g : B → C) (a : A) :
    eitherElim f g (Sum.inl a) = f a := rfl

theorem eitherElim_inr {A B C : Type}
    (f : A → C) (g : B → C) (b : B) :
    eitherElim f g (Sum.inr b) = g b := rfl
\end{minted}
\caption{\texttt{eitherElim}}
\label{lst:ch13_coproduct_universal}
\end{listing}

この \texttt{eitherElim} は、実務のパターンマッチそのものである。左側の処理と右側の処理を渡せば、\texttt{Sum A B} 全体を処理できる。

\begin{listing}[htbp]
\begin{minted}{lean4}
theorem eitherElim_unique_pointwise {A B C : Type}
    (f : A → C) (g : B → C) (h : Sum A B → C)
    (hl : ∀ a, h (Sum.inl a) = f a)
    (hr : ∀ b, h (Sum.inr b) = g b)
    (s : Sum A B) : h s = eitherElim f g s := by
  cases s with
  | inl a =>
      rw [hl a]
      rfl
  | inr b =>
      rw [hr b]
      rfl
\end{minted}
\caption{\texttt{eitherElim\_unique\_pointwise}}
\label{lst:ch13_coproduct_unique}
\end{listing}

積では、同じ入力から二つの出力を作ってペアにした。余積では、二つの入力候補を一つの出力型へ処理する。向きが入れ替わっていることが重要である。

\section{ミニケース：エラーハンドリングとレスポンス型設計}
\label{sec:ch13-api-response}

APIレスポンスを設計するとき、次のような型を考えたくなることがある。

\begin{listing}[htbp]
\begin{minted}{lean4}
structure BadResponse where
  errorCode : Nat
  body      : String
\end{minted}
\caption{\texttt{BadResponse}}
\label{lst:ch13_bad_response}
\end{listing}

この設計は、成功時にもエラーコードがあり、失敗時にも本文があるように読めてしまう。業務仕様が本当に「常に両方ある」なら問題ない。しかし、多くのAPIでは「成功なら本文、失敗ならエラー」である。その場合は余積で表すほうが自然である。

図\ref{fig:ch13-api-response-design}では、エラーと本文を同時に持てる積から、エラーまたは本文だけを持つ余積へ型を変えることで、有効な状態を絞る。

\FigurePlaceholder{fig-ch13-api-response-design.png}{0.90}{レスポンス型の二案}{fig:ch13-api-response-design}

次のLeanコードでは、エラーまたは成功本文を \texttt{Sum} で表す。

\begin{listing}[htbp]
\begin{minted}{lean4}
inductive ApiError where
  | notFound
  | invalidInput

structure ResponseBody where
  status : Nat
  body   : String

abbrev ApiResponse := Sum ApiError ResponseBody

def handleResponse {C : Type}
    (onError : ApiError → C)
    (onOk : ResponseBody → C)
    (r : ApiResponse) : C :=
  eitherElim onError onOk r

theorem handleResponse_ok (onError : ApiError → String)
    (onOk : ResponseBody → String) (body : ResponseBody) :
    handleResponse onError onOk (Sum.inr body) = onOk body := rfl
\end{minted}
\caption{\texttt{ApiResponse} と \texttt{handleResponse}}
\label{lst:ch13_api_response}
\end{listing}

この設計では、レスポンスを使う側が、成功ケースと失敗ケースを明示的に処理する。これにより、「エラーなのに成功値を読もうとした」「成功なのにエラーコードを見てしまった」という混乱を減らせる。

\section{よくある誤解}

\subsection{誤解1：余積は単なる列挙型である}

列挙型は、値を持たない選択肢を表すことが多い。一方、余積は各分岐が値を持てる。\texttt{Sum Error User} なら、左側には \texttt{Error} の値が入り、右側には \texttt{User} の値が入る。したがって、余積は「タグだけ」ではなく「タグ付きデータ」として読むほうが実務に近い。

\subsection{誤解2：エラー型は常に余積にすればよい}

余積は便利だが、すべてのエラー設計を解決するわけではない。たとえば、複数の警告を成功値と一緒に返したい場合は、\texttt{Warnings × Value} のような積が適していることもある。重要なのは、同時に持つべき情報か、排他的な分岐かを見分けることである。

\subsection{誤解3：普遍性は実務には関係ない}

普遍性は抽象的な言葉だが、実務的には「この構造を作る・使う標準的な方法が一つに決まる」という意味を持つ。ペアは射影で特徴づけられ、分岐は各ケースのハンドラで特徴づけられる。この見方は、API設計やDTO設計のレビューで役に立つ。

\begin{warningbox}
積と余積の区別を曖昧にすると、型が業務仕様を正しく表さなくなる。成功値とエラー値を同時に持つ構造にしてしまうと、実装側で「どちらを信じるべきか」を毎回判断する必要が生じる。排他的な状態は、可能な限り余積として表す。
\end{warningbox}

\section{本章のまとめ}

積は「両方持つ」構造であり、ペア、レコード、DTOとして現れる。

余積は「どちらか一方」の構造であり、\texttt{Either}、\texttt{Result}、エラー分岐として現れる。

Leanでは、積を \texttt{Prod} または \texttt{A × B}、余積を \texttt{Sum A B} として扱える。

\texttt{Option A} は、値なしの場合と値ありの場合を持つ小さな余積として読める。

積の普遍性は「二つの射影を持つペア関数が一意に決まる」こととして読める。

余積の普遍性は「左右の処理を与えると、分岐全体の処理が一意に決まる」こととして読める。

APIレスポンスやエラー処理では、同時に持つ情報か排他的な分岐かを区別することが重要である。

\section{次章への接続}

本章では、積と余積によってデータの形を整理した。次章では、\texttt{List.map} や \texttt{Option.map} のように、データの外側の構造を保ったまま中身を変換する考え方へ進む。この考え方が関手である。積や余積で作ったデータ構造を、どのように安全に変換へ持ち上げるかが次の主題になる。
